Replicated data systems consist of multiple agents (or processes) that maintain copies (replicas) of shared data objects. 
These agents perform operations, such as reads and writes, on the replicas. Due to network delays and the asynchronous nature of distributed systems, these operations may not be immediately visible to all agents, leading to potential inconsistencies among replicas. 
The replicated data system abstraction is widely used in various applications, including distributed databases, cloud storage systems, or fully decentralized storage systems such as blockchains or peer-to-peer networks. To some extent, consistency models are also related with weak memory models employed in hardware design~\cite{Burckhardt13}.
An important aspect of replicated data systems is that they often prioritize availability and partition tolerance over strong consistency guarantees. 
Since the strongest consistency model, linearisability~\cite{HerlihyW90}, is too costly to implement in many practical applications, replicated data systems often implement weaker consistency models (often called \emph{isolation level} in the database community). Monotonic reads, monotonic writes, read my writes, eventual consistency, etc are among the most popular ones of the lush jungle of consistency models. Viotti and Vukoli\'{c}~\cite{ViottiVukolic} inventoried 42 consistency models in 2016, and formally established the complex hierarchy they form (see Figure~1 in~\cite{ViottiVukolic}).

Following Burckhardt~\cite{PrinciplesOfEventualConsistency}, a consistency model is a property that an execution of a replicated data system must satisfy. The model of execution of a replicated data system introduced by Burckhardt, called a \emph{history}, bases on a collection of \emph{operations}, each located on a precise process. Each operation also lasts on a time interval that may partially overlap the time interval of another operation located on a different process. 
Like a Gantt diagram, a history therefore combines a partial order on operations and a continuous model of time. The history model is however often too simple for grounding a consistency model. 
Indeed, several consistency models rely on a logical justification of the way how the operations that compose a given history are related to each others. This justification may contradict the partial order and the timings of the operations. Burckhardt~\cite{PrinciplesOfEventualConsistency} therefore also introduced \emph{abstract executions}, that enrich histories with two binary relations among operations, called \emph{visibility} and \emph{arbitration}. 

In this work, we address the problem of automated reasoning on consistency models. Due to the combination of partial orderness and continuous time we just mentioned, it is not obvious which verification models, techniques and tools are the most appropriate for histories and abstract executions (see discussion in conclusion).
Our proposal in this work is to leverage MONA~\cite{MONAguide},  and more generally  MSO-to-automata translations, for automated reasoning on consistency models. MONA models, which are words over a finite alphabet, do not have a lot in common with histories. As a consequence, this requires a significant work for encoding the latter in the former. The choice of MONA also enforces some finiteness limitations, like a bounded, fixed, number of processes, or finitely many values (see conclusion for a detailed discussion on the limitations), which might be compensated by further applications of the automata-based approach (see also discussion in conclusion).

This work stems from an attempt to draw a parallel between communication models and consistency models. This work indeed is partly influenced by previous works of the second author on tree decompositions of message sequence charts for some weakly-synchronous communication models~\cite{DBLP:journals/pacmpl/GiustoFLL23}. Taking a step back in abstraction, we may address the question of automated reasoning on consistency model by transferring the tree decomposition techniques developed for message-passing concurrency to the realm of replicated data systems.

To sum up, we make the following contributions.
\begin{enumerate}
    \item We introduce HistMSO, the monadic second order logic of histories and abstract executions.
    \item We show that HistMSO can express 39 out of the 42 consistency models studied by Viotti and Vukoli\'{c}~\cite{ViottiVukolic}; remarkably, this result is quite straightforward, and follows from a translation of the meta logic used by Viotti and Vukoli\'{c} into HistMSO, based on rather standard MSO encoding techniques for expressing transitive closures, finiteness of sets defined by set comprehension, etc. 
    \item We show that the HistMSO theory of histories is decidable, by reducing the satisfiability problem of HistMSO to the one of MSO over infinite words. Moreover, our reduction also leverages MONA~\cite{MONAguide} as a practical theorem prover for the (weak) HistMSO theory of finite histories.
    \item We introduce $k$-transient visibility, a restriction on the visibility relation, and we show that the HistMSO theory of real-time, $k$-transient abstract executions is decidable.
    \item We show that the partial order defined by the timings of the operations of a given history is the transitive closure of a graph whose cutwidth is bounded by the square of the number of processes; since the treewidth is bounded by the cutwidth, this result induces in particular a tree decomposition of histories bounded by the number of processes and draws a parallel with the second author's previous work.
\end{enumerate}

\paragraph{Outline} Section~\ref{sec:background} recalls the models of histories and abstract executions. In Section~\ref{sec:histmso} we introduce HistMSO logic and show on a few representative examples how to axiomatise the majority of consistency models studied by Viotti and Vukolic. Section~\ref{sec:translation-to-mona} develops the translation of HistMSO to MONA, first for histories, then for $k$-transient real-time abstract executions. Section~\ref{sec:cutwidth} constructs a graph from an history whose transitive closure captures the returns before partial order over operations, and establishes a bound on the cutwidth of the graph that only depends on the number or processes. 
\iflong\else
    Due to space constraints, some details and proofs are omitted and can be found in the full version of the paper~\cite{coget2026histmsologicreasoningconsistency}.
\fi

